\section{Background}
\label{sec:bg-full}

Parallel reasoning and agent programs both expand one committed trunk into several residuals and then keep a single spine~\cite{tot,selfconsistency,react}.
The cost of that expansion is fixed by \emph{when} a system is allowed to reject a child.
Prefix caching can share a block only after the block has been published, disaggregation moves every page the prefiller has already computed, and decode-time pruners score a child only after its prefill and a generated prefix.
\Cref{tab:trad} records the stage.
The rest of this section is why that stage falls after the tokens have been paid for, and what a fan-out looks like once the trunk, the residuals, and the spine are named separately.

\subsection{Parallel Reasoning as a Fan-out}
\label{sec:bg-search}

Chain-of-thought writes one trace~\cite{cot}.
Tree-of-Thoughts, self-consistency, and Monte Carlo tree search write many, then select~\cite{tot,selfconsistency,mctsllm}.
In the Game-of-24 trace of Tree-of-Thoughts, a propose prompt first emits candidate equations and a value prompt then marks each thought sure, maybe, or impossible; breadth-first search keeps a beam of width $b{=}5$~\cite{tot}.
The value judgment reads a thought that has already been generated.
The same order appears in agent programs: a coding or tool-using harness proposes several continuations of one trunk, observes one of them, and aborts the rest~\cite{react,sweagent,openhands}.
On GSM8K and MATH the accuracy gain of that width is real~\cite{gsm8k,math,math500}, and so is the token bill.
EquivPruner and early path pruning start from the same observation: most of the expanded steps do not change the answer that is finally kept~\cite{equivpruner,stoploss}.
What they do not change is the serving order.
The candidate residual is still prefilled, and under disaggregation it is still transferred, before any of those pruners runs.

\subsection{Decode-Time Path Pruning}
\label{sec:bg-prune}

A second line drops a candidate after the model has generated it.
Early-stopping self-consistency draws a window of complete samples and skips later samples only once that window agrees~\cite{esc,selfconsistency}.
Speculative Rejection ranks partial generations at a reward horizon and halts the lower quantile~\cite{specrej}.
Dynamic Parallel Tree Search generates a fixed mini-step on every active node, then early-stops low-confidence children~\cite{dpts}.
STOP organizes these signals by source and learnability and instantiates a learned internal token that fires on a generated prefix~\cite{stoploss}.
EquivPruner instead clusters semantically equivalent actions and keeps one representative of each class~\cite{equivpruner}.
In every case the score is a function of text the model has already produced.

The serving consequence is visible when the three decode-time hosts are run alone on the same GSM8K fan-out (Qwen3-8B, $k{=}16$, $D{=}512$).
ESC, Speculative Rejection, and DPTS each prefill $5008$ tokens.
They differ only after that prefill: ESC decodes to the window ($131072$ tokens, $17.05$\,s), Speculative Rejection cuts at its $256$-token horizon ($98304$ tokens, $13.14$\,s), and DPTS cuts after a $100$-token mini-step ($78336$ tokens, $11.86$\,s).
Any-finished accuracy is $16/16$ on all three.
An earlier decode cut reduces the tokens emitted after the kernel.
It leaves the prefill, and any connector payload computed from that prefill, unchanged.
\S\ref{sec:split} therefore uses a second threshold, read from the residual text before the kernel, and \S\ref{sec:eval-ablate} stacks that gate in front of the same three hosts.

\subsection{Prefix Caching and Disaggregated Transfer}
\label{sec:bg-pd}

Automatic prefix caching and RadixAttention reuse a trunk by hashing a full block, $h=\mathrm{hash}(\mathrm{parent},\;\mathrm{tokens},\;\mathrm{extra})$, after those tokens exist~\cite{vllm,sglang}.
A later request hits only the blocks already in the table.
A same-batch fan-out misses: the trunk has not been published, a partial block is a new block, and the engine materializes $k$ copies of $L$.
The hash is the right index for a replay of the winner.
It is not an admission test for the children that will be aborted.

Prefill--decode disaggregation isolates time-to-first-token from inter-token latency by placing the two phases on separate instances and moving KV across a connector~\cite{distserve,splitwise}.
The controller sends the request to the prefiller and then transfers the pages the prefiller computed.
That path has no admission.
A residual that the harness aborts after \texttt{insert} still occupies the connector, so the transfer equals the prefill (\eqref{eq:xfer-stock}).
Program- and workflow-level schedulers raise the unit of scheduling above the request and still issue the next call as a unique continuation~\cite{autellix,thunderagent,continuum,mori}.
Once KV retention has removed trunk recompute, the remaining time-to-first-token term is the residual suffix of children the harness can name before the observation returns.
\sys reads that residual, aliases the trunk, and places on the connector only the miss tail of the residuals it keeps.

\subsection{Prefill, Decode, and Idle}
Autoregressive inference has two mechanically different phases~\cite{orca,distserve,splitwise}.
\emph{Prefill} consumes the prompt in parallel and is compute-bound; \emph{decode} emits one token (or a small chunk) per step and is memory-bandwidth bound.
PagedAttention~\cite{vllm} stores KV in non-contiguous pages; RadixAttention~\cite{sglang} indexes pages by token prefix so concurrent requests that share a system prompt hit a common trunk.
Two consequences follow for agent workloads.
First, next-turn \ttft{} is the prefill of tokens that miss the cache, plus queuing.
Second, a memory-bound decode batch rarely saturates the tensor cores that a concurrent chunked prefill can use~\cite{sarathi}.
For agents, the prompt of turn $i{+}1$ is a delta on turn $i$, applied after a tool interval of highly variable length~\cite{continuum,mori}.

\subsection{Traditional Methods on a Fan-out}
\label{sec:traditional-app}
We group under \emph{traditional methods} the mechanisms a production stack already applies to one fan-out: prefix caching, hash-indexed prefill, prefill--decode disaggregation, and decode-time pruning.
Each either shares the trunk or rejects a child, and the stage of that operation is fixed.
\Cref{tab:trad} records the stage.
\paragraph{First-expand cost.}
APC indexes a full block by $h=\mathrm{hash}(\mathrm{parent},\;\mathrm{tokens},\;\mathrm{extra})$ after those tokens exist~\cite{vllm}; RadixAttention indexes the same prefix~\cite{sglang}.
A same-batch expansion therefore misses and clones $k$ trunks,
$\Work_{\mathrm{APC}}=kL+\sum_i\ell_i$,
while a later replay of $x_\star$ is a hit ($\Work{=}0$).
Baseline disaggregation transfers the same span, so $\Xfer_{\mathrm{PD}}=\Work_{\mathrm{APC}}$ on the first expand (\eqref{eq:xfer-stock}).
ESC~\cite{esc}, Speculative Rejection~\cite{specrej}, and DPTS~\cite{dpts} evaluate a child only after its prefill and a generated prefix, hence after $\Work$ and, under disaggregation, after $\Xfer$.

\subsection{Fan-out Model}
\label{sec:model}

\Cref{def:step} names the same objects used below.
A parent $u$ holds a committed trunk $x_u$ with $|x_u|{=}L$.
A fan-out is a set of residuals $\{\rho_i\}_{i=1}^{k}$, $|\rho_i|{=}\ell_i$.
The bound winner $\star$ leaves the spine $x_\star=x_u\Vert\rho_\star$, $|x_\star|{=}L{+}\ell_\star$.
Let $b$ be KV bytes per token and $P$ the page size.
The three occupancies are \eqref{eq:mem}:
$M_{\mathrm{clone}}=b(kL+\sum_i\ell_i)$ is the first-expand footprint of recompute, APC, and baseline disaggregation;
$M_{\mathrm{CoW}}=b(L+\sum_i\ell_i)$ is one shared trunk with every residual live;
$M_{\mathrm{spine}}=b(L+\ell_\star)$ is the footprint after abort.
$\Work$ is the token count issued to a prefill kernel before $\star$ is known;
$\Xfer$ is the token count inserted on the P/D connector.
\Cref{tab:notation} lists the remaining symbols.

\begin{table}[t]
\centering
\caption{Notation for Definition~\ref{def:step} and for \eqref{eq:action}--\eqref{eq:xfer}.}
\label{tab:notation}
\scriptsize
\setlength{\tabcolsep}{3.2pt}
\begin{tabular}{@{}llp{0.50\columnwidth}@{}}
\toprule
Symbol & Domain & Meaning \\
\midrule
$L$, $\ell_i$, $\ell_\star$ & $\mathbb{N}$ & trunk; residual $i$; committed residual \\
$\rho_i$, $x_u$ & tokens & residual of child $i$; committed trunk \\
$h_i$ & digest & $\mathrm{hash}(\mathrm{parent},\rho_i[0{:}P),\mathrm{extra})$ \\
$m_i$ & tokens & full-page prefix of $\rho_i$ in $H$, steps of $P$ \\
$s_i$, $s'_i$, $\tau$ & $[0,1]$ & score of $\rho_i$; score of the probe $y_i$; threshold \\
$t_0$, $D$, $d_i$ & tokens & probe length; decode budget; length assigned to $i$ \\
$\alpha$ & $(0,1]$ & early-abort prefix fraction ($0.20$) \\
$p_i$, $q_i$ & $[0,1]$ & admit probability; residual-prefix hit rate \\
$a_i\in\Act$ & $\Act$ & $\mathrm{HS},\mathrm{HP},\mathrm{DS},\mathrm{EA},\mathrm{PF}$ (\eqref{eq:action}) \\
$\Lambda$, $F_i$ & set, $\{0,1\}$ & repeated-loop index set; early-abort predicate \\
$\Keep$ & $\subseteq[k]$ & residuals that remain after admission \\
$\Work_i$, $\Xfer_i$ & tokens & GPU miss tail; connector payload \\
$b$, $P$ & bytes, tok & KV bytes per token; page size \\
$B_t=B^c_t+B^s_t$ & tokens & committed and speculative budget on tick $t$ \\
\bottomrule
\end{tabular}
\end{table}
